\section{Introduction}
\label{sec:introduction}
%版本一
% Automatic speech systems are increasingly moving away from narrowly specialized pipelines toward unified audio-language modeling. In classical ASR, the dominant paradigm evolved from alignment-based and encoder-decoder transduction approaches~\citep{graves2012connectionist,graves2012sequence,chan2015listen} toward systems that combine strong acoustic front ends with large language model decoders~\citep{radford2023robust,malik2021automatic}. In parallel, modern TTS has shifted from hand-engineered pipelines toward generative modeling over increasingly abstract speech representations. A third frontier is now emerging as well: realtime conversational speech agents that must understand paralinguistic signals, respond with low latency, preserve persona, and sound emotionally appropriate within an unfolding interaction. These three trajectories now meet at a common point: speech is no longer treated as a modality that requires a fully separate stack, but as another sequence type that can be mapped into and out of a shared language-centric latent space.

%Despite this convergence, technical descriptions of modern speech systems are still often organized around endpoint tasks. Pretraining, recognition, synthesis, and realtime capability are discussed separately, which is convenient from an engineering perspective but can obscure the actual structure of the model family. In the case of StepAudio 2.5, such a task-wise view is incomplete. The system is more naturally interpreted as one shared audio-language foundation together with three specialization paths: ASR, TTS, and Realtime. Within this family, pretraining establishes the common multimodal prior, ASR shapes that prior into a high-accuracy and low-latency recognizer, TTS adapts the same backbone for controllable expressive generation, and Realtime extends the same foundation toward low-latency human-like spoken dialogue.

% The goal of this report is therefore not merely to juxtapose multiple capabilities within a single document. More importantly, it is to expose the architectural and methodological logic that connects them. The central claim is that StepAudio 2.5 is most naturally understood as an audio-language foundation model whose downstream capabilities differ primarily in data construction, optimization target, decoding regime, and latency constraint, rather than in the underlying multimodal representation itself. Under this view, recognition, synthesis, and realtime spoken interaction are not parallel products loosely assembled around a shared name. They are different operational regimes of one common foundation.

% This perspective has several technical consequences. First, the shared pretraining curriculum should be treated as the central object of analysis, because it establishes the text-audio interface reused by all downstream branches. Second, ASR, TTS, and Realtime can all be interpreted as grounded generation problems, although they differ in what constrains future outputs: acoustic evidence in ASR, semantic control in TTS, and conversational state in Realtime interaction. Third, the most consequential specialization occurs after pretraining. In StepAudio 2.5, ASR advances the quality-efficiency frontier through verifiable multi-token decoding, TTS improves semantic faithfulness and expressivity through context-rich supervision and reinforcement learning, and Realtime introduces persona stability, paralinguistic responsiveness, and turn-level interaction constraints. Finally, the Realtime branch is especially informative because it shows how the same audio-language foundation can support a deployment regime centered on live spoken dialogue.

% These ideas lead to a broader interpretation of the system family. \stepaudio\ is not simply a collection of speech endpoints. It is an instance of a more general thesis: once text and audio are brought into a sufficiently well-shaped shared sequence space, the main task-specific differences increasingly migrate from backbone design to objective design, data design, decoding design, and latency budgeting, even when the family ultimately includes more than one deployed model.

%版本二
% Automatic speech systems are entering a period of architectural convergence, driven by the increasing dominance of large language models: as LLMs became the standard interface for text-based reasoning, treating speech as another sequence type within the same modeling framework became a natural design choice. In ASR, the dominant paradigm has evolved from alignment-based and encoder-decoder transduction approaches~\citep{graves2012connectionist,graves2012sequence,chan2015listen} through large-scale weakly supervised acoustic models such as Whisper ~\citep{radford2023robust}, and more recently toward systems that couple strong acoustic encoders with large language model decoders~\citep{peng2026vibevoiceasrtechnicalreport, an2025funasrtechnicalreport, bai2024seed, shi2026qwen3asrtechnicalreport} . In parallel, TTS has shifted from hand-engineered pipelines toward generative modeling over increasingly abstract speech representations, with commercial systems such as ElevenLabs-v3, Minimax Speech-2.8-hd, and Gemini-Flash-TTS advancing the controllability and expressivity of synthesized speech. A third frontier has emerged in realtime conversational speech agents, exemplified by GPT-realtime, Gemini Live, and Doubao Realtime, that must understand paralinguistic signals, respond with low latency, preserve persona, and remain emotionally appropriate within an unfolding interaction. These three trajectories now meet at a common point: speech is no longer treated as a modality requiring a fully separate stack, but as another sequence type that can be mapped into and out of a shared language-centric latent space, as demonstrated by recent unified audio-language foundations~\citep{wu2025stepaudio2technicalreport, xu2025qwen3omnitechnicalreport}.

% The appeal of this convergence extends beyond consolidating previously separate models into a unified architecture. Traditional cascaded pipelines connect ASR, an intermediate language model, and TTS as isolated stages, inevitably discarding information when speech is reduced to a textual intermediate representation~\citep{tang2024salmonn, borsos2023audiolm, cui2025recent}. A unified audio-language foundation instead preserves speech information end-to-end, allowing paralinguistic cues, emotional state, and conversational context to directly influence recognition, synthesis, and dialogue generation\citep{kim2024paralinguistics, wang2025freezeomni}. Such models also directly leverage the semantic, conversational, and reasoning capabilities already developed in large language models. Under this formulation, audio-language modeling is not merely a matter of replacing task-specific systems with a shared backbone, but of establishing a common representational substrate where information previously lost between stages remains available throughout the interaction process.

% Open-source efforts such as Step-Audio 2 \citep{wu2025stepaudio2technicalreport} and Qwen3-Omni \citep{xu2025qwen3omnitechnicalreport}, alongside large-scale commercial systems including GPT-4o, Gemini~\citep{team2023gemini}, and Doubao, have all moved toward end-to-end audio-language foundations spanning ASR, TTS, and realtime spoken interaction~\citep{yan2025step,cui2025recent,zhang2025mamba,liu2026code}. Despite this shared direction, simultaneously meeting the deployment requirements of all three capabilities within a single model remains challenging. ASR prioritizes accurate and efficient long-form transcription, and TTS emphasizes controllable and expressive synthesis, while realtime interaction further requires low-latency turn-taking together with persona consistency and paralinguistic responsiveness. These objectives are not naturally aligned, and existing unified systems often achieve strong performance on some capabilities while remaining behind specialized systems on others. Closing this gap remains an active focus of audio-language research, without a clearly established solution yet emerging.

% This report presents StepAudio 2.5, building on the Step-Audio line of work \citep{wu2025stepaudio2technicalreport, tian2025step, zhang2026step} to narrow the gap between unified and specialized speech systems. The model adopts a shared audio-language foundation, established through large-scale staged multimodal pretraining that aligns speech and text within a common sequence-modeling interface. On top of this foundation, we view post-training as the primary lever for shaping each capability to its deployment objective, with task-tailored reinforcement learning playing a key role alongside capability-specific supervised fine-tuning and decoding strategies. Rather than treating ASR, TTS, and realtime spoken interaction as separate engineering tracks, we treat them as three specializations of the same audio-language foundation, each refined under its own combination of data, objectives, and post-training signals.

% Concretely, the ASR branch couples the shared decoder with a verifiable multi-token decoding head, exploiting the acoustic determinism of transcription to emit several tokens per step while preserving recognition accuracy under long-form audio. The TTS branch casts speech generation as semantic-to-audio alignment under natural-language control, learned from context-rich supervision over global control descriptions and inline expression directives, and further refined through reinforcement learning over instruction-following and human-preference signals. The Realtime branch is trained for low-latency multi-turn dialogue through a progressive supervised fine-tuning curriculum that introduces persona conditioning and paralinguistic sensitivity, followed by RLHF driven by a generative reward model and explicit interaction rubrics. On standard benchmarks across the three capabilities, StepAudio 2.5 achieves state-of-the-art results in ASR, TTS, and realtime spoken interaction, outperforming both leading unified audio-language models and specialized systems built for each individual task.



% Automatic speech systems are entering a period of architectural convergence, driven by the increasing dominance of large language models: as LLMs became the standard interface for text-based reasoning, treating speech as another sequence type within the same modeling framework became a natural design choice. In ASR, the dominant paradigm has evolved from alignment-based and encoder-decoder transduction approaches~\citep{graves2012connectionist,graves2012sequence,chan2015listen} through large-scale weakly supervised acoustic models such as Whisper~\citep{radford2023robust}, and more recently toward systems that couple strong acoustic encoders with large language model decoders~\citep{peng2026vibevoiceasrtechnicalreport, an2025funasrtechnicalreport, bai2024seed, shi2026qwen3asrtechnicalreport}. In parallel, TTS has shifted from hand-engineered pipelines toward generative modeling over increasingly abstract speech representations, with commercial systems such as ElevenLabs-v3, Minimax Speech-2.8-hd, and Gemini-Flash-TTS advancing the controllability and expressivity of synthesized speech. A third frontier has emerged in realtime conversational speech agents, exemplified by GPT-realtime, Gemini Live, and Doubao Realtime, that must understand paralinguistic signals, respond with low latency, preserve persona, and remain emotionally appropriate within an unfolding interaction. These three trajectories now meet at a common point: speech is no longer treated as a modality requiring a fully separate stack, but as another sequence type that can be mapped into and out of a shared language-centric latent space, as demonstrated by recent unified audio-language foundations~\citep{wu2025stepaudio2technicalreport, xu2025qwen3omnitechnicalreport}.

% The appeal of this convergence extends beyond consolidating previously separate models into a unified architecture. Traditional cascaded pipelines connect ASR, an intermediate language model, and TTS as isolated stages, inevitably discarding information when speech is reduced to a textual intermediate representation~\citep{tang2024salmonn, borsos2023audiolm, cui2025recent}. A unified audio-language foundation instead preserves speech information end-to-end, allowing paralinguistic cues, emotional state, and conversational context to directly influence recognition, synthesis, and dialogue generation~\citep{kim2024paralinguistics, wang2025freezeomni}. Such models also directly leverage the semantic, conversational, and reasoning capabilities already developed in large language models. Under this formulation, audio-language modeling is not merely a matter of replacing task-specific systems with a shared backbone, but of establishing a common representational substrate where information previously lost between stages remains available throughout the interaction process.

% Open-source efforts such as Step-Audio 2 \citep{wu2025stepaudio2technicalreport} and Qwen3-Omni \citep{xu2025qwen3omnitechnicalreport}, alongside large-scale commercial systems including GPT-4o, Gemini~\citep{team2023gemini}, and Doubao, have all moved toward end-to-end audio-language foundations spanning ASR, TTS, and realtime spoken interaction~\citep{yan2025step,wu2025mind,cui2025recent,zhang2025mamba,liu2026code}. Despite this shared direction, simultaneously meeting the deployment requirements of all three capabilities within a single model remains challenging. ASR prioritizes accurate and efficient long-form transcription, and TTS emphasizes controllable and expressive synthesis, while realtime interaction further requires low-latency turn-taking together with persona consistency and paralinguistic responsiveness. These objectives are not naturally aligned, and existing unified systems often achieve strong performance on some capabilities while remaining behind specialized systems on others. Closing this gap remains an active focus of audio-language research.

% This report presents StepAudio 2.5, building on the Step-Audio line of work \citep{wu2025stepaudio2technicalreport, tian2025step, zhang2026step} to narrow the gap between unified and specialized speech systems. The system is most naturally understood not as a collection of parallel endpoints loosely assembled around a shared name, but as a singular audio-language foundation model guided by a central thesis:

% \begin{tcolorbox}[colback=gray!10, colframe=gray!10, boxrule=0pt, arc=2mm, halign=center]
% \textit{Once text and audio share a well-shaped representational space, the differences among downstream tasks migrate away from architecture toward operational regimes: data, objectives, and decoding constraints.}
% \end{tcolorbox}

% Following this insight, we view post-training as the primary lever for shaping each capability to its specific deployment objective. Rather than treating ASR, TTS, and realtime interaction as separate engineering tracks, we refine the shared multimodal prior through task-tailored reinforcement learning, capability-specific supervised fine-tuning (SFT), and specialized decoding strategies. Concretely, the ASR branch advances the quality-efficiency frontier by coupling the shared decoder with a verifiable multi-token decoding head, exploiting acoustic determinism to emit multiple tokens per step. The TTS branch adapts the backbone for controllable generation via semantic-to-audio alignment, integrating context-rich supervision with reinforcement learning over instruction-following and human-preference signals. Finally, the Realtime branch extends the foundation toward low-latency spoken dialogue through progressive SFT for persona and paralinguistic sensitivity, followed by RLHF driven by a generative reward model and explicit interaction rubrics. On standard benchmarks across the three capabilities, StepAudio 2.5 achieves state-of-the-art results, outperforming both leading unified audio-language models and specialized systems built for individual tasks.

Automatic speech systems are entering a period of architectural convergence, driven by the increasing dominance of large language models (LLMs): as LLMs became the standard interface for text-based reasoning, treating speech as another sequence type within the same modeling framework became a natural design choice. In automatic speech recognition (ASR), the dominant paradigm has evolved from alignment-based and encoder-decoder transduction approaches~\citep{graves2012connectionist,graves2012sequence,chan2015listen} through large-scale weakly supervised acoustic models such as Whisper~\citep{radford2023robust}, and more recently toward systems that couple strong acoustic encoders with LLM decoders~\citep{peng2026vibevoiceasrtechnicalreport, an2025funasrtechnicalreport, bai2024seed, shi2026qwen3asrtechnicalreport}. In parallel, text-to-speech (TTS) synthesis has shifted from hand-engineered pipelines toward generative modeling over increasingly abstract speech representations, with commercial systems such as ElevenLabs-v3, Minimax Speech-2.8-hd, and Gemini-Flash-TTS advancing the controllability and expressivity of synthesized speech. A third frontier has emerged in realtime conversational speech agents, exemplified by GPT-realtime, Gemini Live, and Doubao Realtime, that must understand paralinguistic signals, respond with low latency, preserve persona, and remain emotionally appropriate within an unfolding interaction. These three trajectories now meet at a common point: speech is no longer treated as a modality requiring a fully separate stack, but as another sequence type that can be mapped into and out of a shared language-centric latent space, as demonstrated by recent unified audio-language foundations~\citep{wu2025stepaudio2technicalreport, xu2025qwen3omnitechnicalreport,liu2026omni}.

The appeal of this convergence extends beyond consolidating previously separate models into a unified architecture. Traditional cascaded pipelines connect ASR, an intermediate language model, and TTS as isolated stages, inevitably discarding information when speech is reduced to a textual intermediate representation~\citep{tang2024salmonn, borsos2023audiolm, cui2025recent}. A unified audio-language foundation instead preserves speech information end-to-end, allowing paralinguistic cues, emotional state, and conversational context to directly influence recognition, synthesis, and dialogue generation~\citep{kim2024paralinguistics, wang2025freezeomni,li2026depflow,xuan2023new,deng2025multi}. Such models also directly leverage the semantic, conversational, and reasoning capabilities already developed in LLMs. Under this formulation, audio-language modeling is not merely a matter of replacing task-specific systems with a shared backbone, but of establishing a common representational substrate where information previously lost between stages remains available throughout the interaction process.

Open-source efforts such as Step-Audio 2 \citep{wu2025stepaudio2technicalreport} and Qwen3-Omni \citep{xu2025qwen3omnitechnicalreport}, alongside large-scale commercial systems including GPT-4o, Gemini~\citep{team2023gemini}, and Doubao, have all moved toward end-to-end audio-language foundations spanning ASR, TTS, and realtime spoken interaction~\citep{yan2025step,wu2025mind,wu2025chronological,zhang2026sla,cui2025recent,zhang2025mamba,liu2026code}. Despite this shared direction, simultaneously meeting the deployment requirements of all three capabilities within a single model remains challenging. ASR prioritizes accurate and efficient long-form transcription, and TTS emphasizes controllable and expressive synthesis, while realtime interaction further requires low-latency turn-taking together with persona consistency and paralinguistic responsiveness. These objectives are not naturally aligned, and existing unified systems often achieve strong performance on some capabilities while remaining behind specialized systems on others. Closing this gap remains an active focus of audio-language research.

This report presents StepAudio 2.5, building on the Step-Audio line of work \citep{wu2025stepaudio2technicalreport, tian2025step, zhang2026step} to narrow the gap between unified and specialized speech systems. The system is most naturally understood not as a collection of parallel endpoints loosely assembled around a shared name, but as a singular audio-language foundation model guided by a central thesis:

\begin{tcolorbox}[colback=gray!10, colframe=gray!10, boxrule=0pt, arc=2mm, halign=center]
\textit{Once text and audio share a well-shaped representational space, the differences among downstream tasks migrate away from architecture toward operational regimes: data, objectives, and decoding constraints.}
\end{tcolorbox}

Following this insight, we view post-training as the primary lever for shaping each capability to its specific deployment objective. Rather than treating ASR, TTS, and realtime interaction as separate engineering tracks, we refine the shared multimodal prior through a unified alignment paradigm. Crucially, we move beyond basic supervised fine-tuning (SFT) by establishing Reinforcement Learning from Human Feedback (RLHF) as the central mechanism for capturing nuanced human preferences and paralinguistic behaviors. We complement this RLHF-driven alignment with capability-specific SFT and specialized decoding strategies. Concretely, the ASR branch advances the quality-efficiency frontier by coupling the shared decoder with a verifiable multi-token decoding head, exploiting acoustic determinism to emit multiple tokens per step. The TTS branch adapts the backbone for controllable generation via semantic-to-audio alignment, integrating context-rich supervision with human-preference-driven RLHF. Finally, the Realtime branch extends the foundation toward low-latency spoken dialogue through progressive SFT for persona and paralinguistic sensitivity, followed by RLHF driven by a generative reward model and explicit interaction rubrics. On standard benchmarks across the three capabilities, StepAudio 2.5 achieves state-of-the-art results, outperforming both leading unified audio-language models and specialized systems built for individual tasks.